% Generated draft: requires mathematical and source review.
\documentclass[11pt]{article}
\usepackage[margin=25mm]{geometry}
\usepackage{fontspec}
\setmainfont{Latin Modern Roman}
\usepackage{amsmath,amssymb,mathtools,unicode-math}
\setmathfont{Latin Modern Math}
\usepackage{enumitem,needspace,xurl,hyperref}
\setlength{\parindent}{0pt}
\setlength{\parskip}{5pt}
\setlength{\emergencystretch}{3em}
\newcommand{\N}{\mathbb{N}}
\newcommand{\Z}{\mathbb{Z}}
\newcommand{\Q}{\mathbb{Q}}
\newcommand{\R}{\mathbb{R}}
\newcommand{\C}{\mathbb{C}}
\newcommand{\F}{\mathbb{F}}
\newcommand{\A}{\mathbb{A}}
\newcommand{\PP}{\mathbb{P}}
\newcommand{\E}{\mathbb{E}}
\newcommand{\eps}{\varepsilon}
\newcommand{\dd}{\,\mathrm{d}}
\begin{document}
\section*{UnsolvedMath ID: 801065}
\textbf{Source problem code:} catalogue entry
\subsection{Definitions and mathematical statement}
Let \(\mathbb{T}=\mathbb{R}/\mathbb{Z}\), equipped with normalized Lebesgue measure \(m\), and let \(\mathbb{T}_{\mathrm{irr}}=\{x+\mathbb{Z}:x\in\mathbb{R}\setminus\mathbb{Q}\}\). The supplied question begins with the qualification ``If the answer is yes,'' referring to an antecedent question absent from the record. Subject to that unresolved qualification, retain the source's notation
\[
\Theta=\{\theta\in\mathbb{T}_{\mathrm{irr}}:H_2\text{ has a Herman ring with rotation number }\theta\}.
\]
The supplied record does not define \(H_2\), specify whether it denotes a map or a class of maps, or define the intended notion of a Herman ring. Thus the membership condition above remains source-dependent and does not yet specify a fully determined set.

\textbf{Question (Berger--Bedford, Question 17).} With those source-dependent terms resolved and the preceding affirmative-answer qualification satisfied, does
\[
m(\Theta)>0
\]
hold? More specifically, is the class \(0+\mathbb{Z}\) a Lebesgue density point of \(\Theta\)? Explicitly, writing
\[
I_r=\{x+\mathbb{Z}:x\in(-r,r)\},\qquad 0<r<\tfrac12,
\]
the latter question asks whether
\[
\lim_{r\downarrow0}\frac{m(\Theta\cap I_r)}{2r}=1.
\]
These measure formulations presuppose Lebesgue measurability of \(\Theta\), which the supplied record does not establish. Density is taken in the ambient circle \(\mathbb{T}\); it does not require \(0+\mathbb{Z}\) to belong to \(\Theta\). The density assertion is stronger than positive measure, rather than an equivalent reformulation.
\subsection{Short English statement}
Conditionally on an unspecified preceding question having an affirmative answer, Berger and Bedford ask whether the irrational rotation numbers realized by Herman rings for the source's undefined \(H_2\) form a set of positive Lebesgue measure, and more specifically whether that set has density one at zero. The missing antecedent and definition of \(H_2\) require clarification.
\subsection{Sources}
\begin{itemize}
\item Supplied local UnsolvedMath registry record, permanent ID \(801065\), exact source problem code \texttt{catalogue entry}, titled ``Henon Maps: Open Problems --- Question 17.'' The registry attributes the problem to Berger and Bedford, with proposed year 2023, and labels it open; this status is supplied metadata, not an independently verified mathematical or literature assessment.
\item Berger, Bedford, \emph{H\'enon Maps: A List of Open Problems} (2023), Question 17. Supplied source URL: \url{https://arxiv.org/abs/2312.03907}. The provenance records extraction by \texttt{arxiv\_latex} and access on 2026-10-16. This access date and attribution are supplied provenance and have not been independently verified.
\end{itemize}
\end{document}
